\ficha{Summerhill (2015), Inglorious Revolution}{summerhill2015}

\begin{identificacao}
SUMMERHILL, William R. \textit{Inglorious Revolution: political institutions, sovereign debt, and financial underdevelopment in Imperial Brazil}. New Haven: Yale University Press, 2015. 349 p. (Yale Series in Economic and Financial History).

Livro de história econômica, em inglês. Leitura seletiva: introdução (p. 1-18), capítulo 2, ``Sovereign Borrowing and Imperial Debt Policy'' (p. 19-44), capítulo 3, ``Tropical Credibility on Lombard Street'' (p. 45-79), e capítulo 8, ``Fall from Grace'' (p. 215-228), na íntegra; capítulo 5, ``Turning Points'' (p. 121-150), por alto, apenas para registrar a tese. Não foram lidos os capítulos 4 (dívida interna), 6 (incorporação) e 7 (bancos comerciais), nem os três apêndices. Edição publicada, sem versão preliminar.
\end{identificacao}

\bloco{Problema e tese}
O livro pergunta como o Império conseguiu tomar emprestado durante sessenta e cinco anos, em Londres e no Rio, sem repúdio nem default, e ao mesmo tempo deixou de produzir o desenvolvimento financeiro privado que a literatura associa a esse feito. A resposta tem duas metades ligadas pela mesma causa. A Constituição de 1824 transferiu ao parlamento, e sobretudo à câmara dos deputados, a autoridade de tributar, gastar e contrair dívida, o que amarrou as mãos do monarca e tornou crível o compromisso de pagar. As mesmas instituições, ao concentrar toda a política financeira no governo central e responder a um eleitorado censitário estreito, restringiram a incorporação e a entrada bancária. O Brasil imperial é, por isso, um contraexemplo da proposição de North e Weingast segundo a qual o compromisso crível com a dívida soberana leva necessariamente ao desenvolvimento financeiro (p. 5).

\bloco{Argumento}
O livro inteiro sustenta que a ruptura constitucional de 1822 e 1824 foi, em termos financeiros, uma revolução ``inglória'' (p. 11): copiou da Inglaterra pós-1688 a parte que garante o credor do Estado e não copiou a parte que abre o crédito ao empresário. Os oito capítulos testam essa assimetria em quatro frentes, a saber, como o Estado conseguiu tomar emprestado, que papel tiveram os mercados interno e externo, o que moveu o risco de default depois de obtido o acesso, e qual a relação entre as instituições que sustentaram a dívida soberana e os obstáculos ao financiamento privado (p. 13). Do recorte lido saem os passos seguintes.

Primeiro, o problema é político, não fiscal. \chave{Sovereign debt subordinates economics to politics}{p. 19}. Um terço dos defaults desde 1820 ocorreu em bons tempos econômicos, o que desloca a explicação da capacidade de pagar para a disposição de pagar (p. 21). O modelo do capítulo 2 deriva daí que penalidade forte, monitoramento e controle fiscal por parte dos credores baixam o custo e elevam o teto de endividamento.

Segundo, o desenho brasileiro tinha as três peças. A câmara detinha iniciativa em tributos e despesas, e podia derrubar o gabinete; o ministro da Fazenda respondia criminalmente; a abdicação de 1831 provou que nem o imperador era imune (p. 26). O teste qualitativo é a derrota, por mais de dois para um, da proposta de moratória externa apresentada pelo ministro José Inácio Borges em junho de 1831, quando ficou claro que credores internos e externos tinham interesse comum (p. 29). O monitoramento entre sessões coube à Junta Administrativa da Caixa de Amortização, criada em 1827 com cinco ``capitalistas nacionais'' entre os maiores portadores de apólices, com acesso contínuo às contas do Tesouro (p. 33).

Terceiro, a sustentabilidade fiscal é testada por econometria. A série do déficit total é estacionária, o saldo primário e a dívida são cointegrados, e a regressão em primeiras diferenças mostra resposta positiva do primário à dívida: o governo elevava o superávit exigido em cerca de 8 a 11 libras para cada 100 libras adicionais de dívida (p. 42). Em apenas onze de sessenta e oito anos houve superávit total, mas o saldo primário era em geral positivo, o que desmente a tese de déficit crônico repetida por Caio Prado Júnior e Celso Furtado (p. 7 e p. 37).

Quarto, a credibilidade em Londres é medida por três indicadores construídos a partir dos contratos de empréstimo, e não por narrativa. O custo ex ante de capital de cada empréstimo é a taxa interna de retorno que iguala o dinheiro efetivamente recebido ao fluxo descontado de juros, amortização e comissões do banco (p. 71). O prêmio de risco é a diferença entre esse custo e o rendimento dos consols britânicos na semana da assinatura do contrato (p. 72). A probabilidade de default é extraída do preço de emissão sob dois supostos, neutralidade ao risco e distribuição de dois pontos entre cumprimento integral e repúdio total (p. 74). Os três convergem: o custo cai de 13,89\% em 1829 para 5,12\% em 1889, o prêmio de risco de 10,41 para 2,47 pontos, e a probabilidade de repúdio de 7,6\% para 1,9\% (p. 72 e p. 76). O prazo é o quarto indicador, qualitativo: nenhum empréstimo londrino teve prazo inferior a vinte anos, e o de conversão de 1889 venceria em cinquenta e seis (p. 66). Em uma era de consols entre 2,45\% e 4\%, o risco-país brasileiro é visível, mas cadente (p. 73).

Quinto, o capítulo 5, lido por alto, mede o risco no mercado secundário com dados semanais dos dois lados do Atlântico e localiza quebras estruturais. A tese que interessa é negativa: rejeita-se a hipótese de que a reputação de pagamento pontual explique a queda do prêmio; os mercados olhavam para a frente, para eventos políticos e guerras, não para o histórico (p. 148). Londres reagia a conflito externo, o Rio a conflito interno e ao risco de inflação, porque o portador de apólice corria risco cambial que o portador de bônus em libras não corria (p. 124). O prêmio médio foi de cerca de 340 pontos-base em Londres e 370 no Rio (p. 149).

\bloco{Conceitos e definições operacionais}
Compromisso crível é definido de modo institucional e mensurável: existe quando credores ou seus representantes detêm direito de monitoramento e de controle fiscal encravado em regra constitucional cuja violação gera crise política, não apenas embargo de crédito (p. 24). Credibilidade não é atributo, é preço, e se lê no prêmio de risco e no teto de endividamento, este identificado com o estoque de dívida externa após cada novo empréstimo (p. 67).

Padrão monetário: o mil-réis ``raramente'' foi conversível em ouro a taxa fixa e flutuou na maior parte do século XIX (p. xiii). A lei de 1846 fixou a paridade em 27 pence por mil-réis, mas Summerhill trata essa paridade como nocional, ``um alvo de taxa de câmbio que não era sustentado de modo continuado'' (p. xiii). O papel do Tesouro não tinha lastro metálico e rendia senhoriagem (p. 84). A lei bancária de 1860 só permitia emissão com conversibilidade em ouro na razão de um para um (p. 191). Ouro Preto autorizou em julho de 1889 emissão conversível de até três vezes o capital em espécie do banco (p. 209).

Pecado original financeiro: dívida denominada em moeda estrangeira. O Império escapou da forma mais forte porque emitiu longo, em casa, na própria moeda inconversível, e por isso combinou dívida indexada e dívida nominal, o que reduzia seu incentivo a inflacionar (p. 79).

\bloco{Evidência e método}
Estudo de caso quantitativo de longo prazo, 1822 a 1889, com prolongamento até 1898 e menções até 1982. Fontes: contratos originais dos dezoito empréstimos externos, arquivo Rothschild em Londres, Arquivo Nacional, biblioteca e arquivo do Ministério da Fazenda, relatórios do ministro da Fazenda, Coleção das Leis do Império, almanaques do Rio, imprensa londrina, cotações semanais de bônus em Londres e no Rio. Técnicas: séries originais de receita, despesa, dívida e serviço; testes de raiz unitária, cointegração e correção de erros; regressão em primeiras diferenças com controles de ciclo à Barro; cálculo de taxa interna de retorno por contrato; teste de quebras estruturais múltiplas nas séries de prêmio de risco. Contrafactual explícito: menos centralização política teria antecipado a liberalização da incorporação e ampliado a intermediação financeira (p. 9).

\bloco{Agentes e vertentes}
Referências teóricas: North e Weingast como alvo, com Robinson, Salvucci, Reinhart e Rogoff, e a literatura sobre cidades-estado italianas e República Holandesa. Agentes brasileiros: Pedro I, José Inácio Borges, Cunha Mattos, Bernardo de Vasconcelos, Montezuma, Miguel Calmon du Pin e Almeida, visconde de Itaboraí, visconde do Rio Branco, Silveira Martins, Francisco Belisário Soares de Souza, João Alfredo, visconde de Ouro Preto, Rui Barbosa, Deodoro da Fonseca, Floriano Peixoto e Campos Salles. Banqueiros: Nathan Mayer Rothschild e N. M. Rothschild \& Sons, agente financeiro oficial desde 1855, responsável por catorze das dezoito emissões europeias (p. 50), além de Samuel \& Phillips, Goldsmid, Erlanger, Hambro, Baring e Banque de Paris et Pays Bas. Historiografia contestada: Caio Prado Júnior, Celso Furtado e José Honório Rodrigues, pela tese do déficit crônico (p. 7).

\bloco{Críticas e limites}
O livro termina em 1889 e trata todo o período republicano em um capítulo final de quatorze páginas. A Caixa de Conversão não aparece; a adoção do padrão-ouro em 1906 recebe um parágrafo (p. 227). Não há análise do desenho da Caixa, do nível da taxa, do limite de emissão ou do debate público. O próprio autor admite que o efeito das restrições sobre o financiamento privado é impossível de calcular com precisão quantitativa (p. 9) e que a probabilidade de default é enviesada pelos supostos de repúdio total e neutralidade ao risco (p. 75). O eixo do livro é dívida e incorporação, não padrão monetário: metalismo, papelismo e o debate doutrinário do Império estão ausentes.

\begin{dialogo}
\item Amarra institucional da Caixa. Procurar nos jornais e nos Documentos Parlamentares os argumentos sobre a competência da Caixa frente ao Tesouro e ao Banco do Brasil, e sobre a custódia do ouro em 1911, com vocabulário de época como ``fundo de garantia'', ``fundo de resgate'' e ``lastro inviolável''. Testa se os agentes de 1906 raciocinavam como Summerhill descreve a Junta da Caixa de Amortização (p. 33), isto é, se a credibilidade era buscada em regra que retirasse discricionariedade do Executivo, e não em promessa de política.
\item O banqueiro como certificador. Procurar as menções aos Rothschild, à agência de Londres e ao funding loan de 1898 em O Paiz e no Correio da Manhã. Testa a hipótese de Summerhill de que a assinatura da casa emissora sinalizava qualidade ao mercado (p. 51), e se a imprensa brasileira via o banqueiro como aval ou como tutela.
\item A paridade de 27 pence. Procurar se os defensores do par legal de 27 pence invocavam a lei de 1846 e a conversibilidade alcançada por Ouro Preto em 1889, e se os defensores de 12 ou 15 pence tratavam essa paridade como nocional. Testa a leitura de Summerhill de que 27 pence sempre foi alvo, nunca regime sustentado (p. xiii e p. 78).
\item Padrão-ouro como sinal a credores externos. Procurar a justificativa da Caixa pela recuperação do crédito externo perdido nos anos 1890 e pela facilitação de novos empréstimos. Summerhill afirma que a adoção de 1906 foi \chave{prompted as much by desperation as by anything else}{p. 227} e funcionou por poucos anos. Testa se o argumento do crédito externo aparece na imprensa como motivo declarado ou apenas como efeito esperado.
\item Crédito soberano e crédito privado. Procurar se os jornais ligam a Caixa ao crédito para a lavoura, para a indústria e para os bancos, ou se tratam os dois planos como separados. Testa, para 1906 a 1914, o desacoplamento que Summerhill documenta para o Império entre custo de captação do Estado e taxas do setor privado (p. 3).
\end{dialogo}

\usonocapitulo{Parte II, como caso central de compromisso crível construído por instituição política e não por regra monetária, e como fonte das medidas de credibilidade que o capítulo pode reutilizar, a saber, prêmio sobre consols, custo ex ante de capital, prazo dos títulos e probabilidade implícita de default. Sustenta duas afirmações do capítulo: que credibilidade soberana e desenvolvimento financeiro privado são coisas distintas, e que os mercados precificavam expectativa política, não reputação acumulada. Serve também de ponte para a parte III, ao estabelecer a herança institucional e cambial que os agentes de 1906 recebem, isto é, uma paridade legal de 27 pence nunca sustentada, uma conversibilidade brevemente alcançada em 1888 e 1889 e perdida na primeira década republicana, e um crédito externo arruinado pelo default de 1898. Para o desenho da Caixa em si o livro não serve, e isso deve ser dito no capítulo, porque Summerhill não a analisa.}
